\chapter{System Design and FPGA Accelerator}
\label{cha:design}

This chapter describes the proposed methodology. Section~\ref{sec:background} first
introduces the two architectural notions on which the rest of the chapter relies: the
SwiGLU feed-forward block and the LFM2 hybrid backbone. Section~\ref{subsec:cpuexec}
then identifies the dominant computational kernel through fine-grained profiling,
Section~\ref{sec:orchestration} details the CPU--FPGA orchestration that offloads that
kernel, and Section~\ref{sec:fpga} presents the design of the FPGA accelerator itself.

\section{Background}
\label{sec:background}

\subsection{The Feed-Forward Block and SwiGLU}
\label{subsec:swiglu_bg}
Modern transformer language models alternate two kinds of sub-layer within each block:
a token-mixing operator (self-attention or, in LFM2, a short convolution) and a
position-wise feed-forward network (FFN). The FFN is applied independently to every
token's hidden vector and is responsible for most of a model's parameters and
arithmetic. A classic FFN expands the $d$-dimensional hidden vector to a larger
intermediate dimension $d_{\text{ff}}$ through one linear projection, applies a
pointwise nonlinearity, and projects back to $d$ through a second projection. Gated
Linear Unit (GLU) variants replace the single expansion with two parallel projections
combined by an element-wise gate, which improves quality at an equal parameter budget
and has become standard in recent large language models. LFM2 uses the SwiGLU variant.

SwiGLU combines a SiLU-gated branch with a linear branch. The SiLU (also called Swish)
activation is $\text{SiLU}(z) = z\,\sigma(z)$, where $\sigma$ is the logistic sigmoid;
it is a smooth, non-monotonic function that acts as a soft gate, passing large positive
activations while suppressing negative ones. Given the input activation $x \in
\mathbb{R}^{d}$ (here $d = 2048$), SwiGLU computes a gate projection $X_1 = x W^{T}$ and
an up projection $X_2 = x V^{T}$, both mapping to the intermediate dimension
$d_{\text{ff}}$ (here $8192$); it applies SiLU to the gate, multiplies the two branches
element-wise, and projects the result back to $d$ with the down projection $W_{down}$:
\begin{equation}
\text{FFN}(x) = (\text{SiLU}(x W^{T}) \odot (x V^{T})) W_{down}^{T}
\label{eq:ffn}
\end{equation}
where $\odot$ denotes element-wise multiplication. LFM2 applies no bias terms in the
FFN, which simplifies the arithmetic and is exploited by the hardware design in
Section~\ref{sec:fpga}.

The cost of this expressiveness is that SwiGLU uses three large weight matrices ($W$ and
$V$ of size $d \times d_{\text{ff}}$ and $W_{down}$ of size $d_{\text{ff}} \times d$)
rather than the two of a classic FFN. At $d = 2048$ and $d_{\text{ff}} = 8192$ each
matrix holds $2048 \times 8192 \approx 16.8$\,M weights, so a single block's FFN
contains roughly $50$\,M parameters and dominates both the storage footprint and the
multiply-accumulate count of the block. This is what makes the feed-forward path the
natural acceleration target identified in Section~\ref{subsec:cpuexec}.

\subsection{The LFM2 Hybrid Backbone}
\label{subsec:lfm2_bg}
LFM2 is a family of edge-first language models from Liquid AI, designed through a
hardware-in-the-loop architecture search that targets CPU latency and memory directly
rather than parameter count alone~\cite{amini2025lfm2technicalreport}. The resulting
backbone is a minimal hybrid: most blocks use an inexpensive gated short convolution
for token mixing, and only a small minority use grouped-query attention
(GQA)~\cite{ainslie2023gqa} to recover long-range modeling. This combination delivers up to roughly $2\times$ faster
prefill and decode on CPUs than similarly sized attention-only models, which is what
makes LFM2 a credible starting point for edge deployment.

The 1.2-billion-parameter backbone used in this work has the hyperparameters listed in
Table~\ref{tab:lfm2_dims}. It comprises 16 transformer blocks over a hidden dimension
of 2048, of which 10 are gated short-convolution blocks and 6 are GQA blocks; every
block, regardless of its mixer, contains a SwiGLU FFN that expands to a feed-forward
dimension of 8192.

\begin{table}[htbp]
\centering
\caption{Hyperparameters of the LFM2-1.2B language backbone.}
\label{tab:lfm2_dims}
\begin{tabular}{lr}
\hline
\textbf{Hyperparameter} & \textbf{Value} \\ \hline
Total parameters & 1.2\,B \\
Transformer blocks & 16 \\
Hidden dimension ($d_{\text{model}}$) & 2048 \\
Feed-forward dimension ($d_{\text{ff}}$) & 8192 \\
Gated short-convolution blocks & 10 \\
Grouped-query attention blocks & 6 \\
Attention heads / KV groups / head dim & 32 / 8 / 64 \\
Convolution kernel size ($k$) & 3 \\
Normalization & pre-norm RMSNorm \\
Positional encoding & RoPE (attention blocks) \\ \hline
\end{tabular}
\end{table}

Internally, each block follows the pre-norm residual structure illustrated in
Figure~\ref{fig:lfm2_block}: the hidden state is normalized with RMSNorm, passed
through the sequence mixer, and added back to the residual stream; it is then
normalized again, passed through the SwiGLU FFN, and added back. The sequence mixer is
either a gated short convolution (a depthwise causal convolution of kernel size
$k = 3$ surrounded by input-dependent multiplicative gating, which provides cheap local
mixing with favourable cache behaviour on CPUs) or grouped-query attention, which
shares key and value projections across groups of query heads (here 32 query heads over
8 key/value groups) to reduce key--value cache traffic, augmented with rotary position
embeddings and query--key normalization.

\begin{figure}[htb!]
\centering
\includegraphics[width=0.9\linewidth]{lfm2_architecture.png}
\caption{LFM2 architecture, reproduced from the LFM2 technical
report~\cite{amini2025lfm2technicalreport}. Each block alternates a sequence-mixing
operator (a gated short convolution or grouped-query attention) with a SwiGLU
feed-forward block, both wrapped in pre-norm residual connections. The architecture
supports both dense and mixture-of-experts (MoE) variants, and the MoE experts are
themselves SwiGLU blocks. This work uses the dense 1.2B variant and offloads the SwiGLU
feed-forward block to the FPGA.}
\label{fig:lfm2_block}
\end{figure}

Two properties of this structure shape the rest of the chapter. First, because a SwiGLU
FFN appears in all 16 blocks, the feed-forward path aggregates to the single largest
share of per-token computation, as the profiling in Section~\ref{subsec:cpuexec}
confirms. Second, the FFNs are never adjacent: consecutive feed-forward blocks are
always separated by a full sequence mixer and two normalization and residual steps. In
the heterogeneous design of Section~\ref{sec:orchestration}, those intervening operators
remain on the CPU, so the accelerator cannot chain one FFN directly into the next;
instead, every generated token incurs sixteen separate CPU--FPGA handoffs, one per
block.

\section{Profiling under CPU Execution}
\label{subsec:cpuexec}

To address the research question of enabling conversational-quality inference within
the constraints of edge devices, this work adopts a systems-driven, top-down
methodology that combines model profiling with targeted hardware acceleration. Rather
than attempting to accelerate the entire conversational pipeline, the approach focuses
on identifying the computational bottlenecks of the language backbone that powers the
selected speech model and mapping the most demanding operations to dedicated hardware.

The starting point of the approach is the LFM2-Audio architecture, whose
conversational capabilities are supported by a 1.2\,B-parameter language backbone. To
understand its execution characteristics on embedded hardware, the inference engine
implemented by llama.cpp was modified to instrument the execution of individual
architectural blocks. This profiling allows the collection of fine-grained timing
information for each major component of the model during inference.

\begin{table}[htbp]
\centering
\caption{LFM2.5-1.2B CPU execution profile.}
\label{tab:cpu_profile}
\begin{tabular}{lrr}
\hline
\textbf{Architectural Block} & \textbf{Time ($\mu$s)} & \textbf{\% of Total} \\ \hline
MatMul + SwiGLU & 10,136,836 & 68.08\% \\
Final Linear Projection & 1,631,552 & 10.96\% \\
Gated Short Convolution & 1,555,396 & 10.44\% \\
GQA Attention (MatMul) & 1,318,155 & 8.85\% \\
Unmapped Nodes & 174,103 & 1.17\% \\
Memory Ops \& KV Cache & 28,330 & 0.19\% \\
RMSNorm / LayerNorm & 27,328 & 0.18\% \\
Residual Connections \& Addition & 11,446 & 0.08\% \\
Input Embedding Lookup & 4,006 & 0.03\% \\
Positional Encodings (RoPE) & 1,840 & 0.01\% \\ \hline
\end{tabular}
\end{table}

Table~\ref{tab:cpu_profile} summarizes the resulting execution profile on a quad-core
Arm\textregistered{} Cortex\textregistered{}-A53 platform. These results indicate that
the performance of the model is primarily determined by a small set of dense linear
algebra operations. In particular, the feed-forward MatMul + SwiGLU blocks represent
the principal computational bottleneck, followed by the projection layers and
attention-related matrix multiplications. Since these operations exhibit regular and
predictable memory access patterns, they are well suited for hardware acceleration via
spatial architectures that can exploit custom dataflow and data reuse.

Amdahl's Law~\cite{Amdahl1967ValidityOT} bounds the theoretical benefit of this
strategy: with the feed-forward MatMul + SwiGLU block (Eq.~\eqref{eq:ffn}) accounting
for approximately 68\% of total execution time, the maximum achievable system speedup
is limited to $1/(1-0.68) \approx 3.1\times$, even if the accelerated portion were to
execute instantaneously. In practice, the non-FFN overhead was independently measured
at approximately 126\,ms per token by substituting a pass-through IP for the
accelerator, confirming that the remaining operations (attention, normalization, and
residual additions) set a hard throughput ceiling near 8 tokens per second on this
platform.

Based on this analysis, the proposed system focuses on offloading the most
computationally intensive component of the model to an FPGA-based accelerator while
leaving control logic and lightweight operations on the CPU. This selective
acceleration strategy enables significant reduction of inference latency at low thread
counts without requiring full hardware implementation of the entire model.

\section{Orchestration of the Computation Handoff}
\label{sec:orchestration}

Following the identification of the dominant computational bottleneck, the llama.cpp
inference engine was extended with a hardware offloading mechanism that intercepts the
execution of the feed-forward MatMul + SwiGLU block and orchestrates a handoff to a
custom accelerator implemented in the FPGA fabric. Inference is executed using a
modified \texttt{llama.cpp} CPU backend. Model weights are loaded from GGUF, and for
each token a GGML computation graph is built (attention, norms, FFN). The graph is
executed node-by-node via \texttt{ggml\_graph\_compute}, which dispatches each
operation to CPU kernels by default.

To accelerate the dominant FFN block, the LFM2 builder replaces the standard FFN
sequence with a fused operator (\texttt{ggml\_swiglu\_fused\_hw}) when the environment
variable \texttt{LLAMA\_SWIHW=1} and the tensor shapes and types match the accelerator
interface. The layer index is stored in the node to retrieve the correct weights. At
execution time, this fused node triggers a custom kernel that manages CPU--FPGA
interaction. Weights are cached once per layer in a physically contiguous,
DMA-coherent buffer (udmabuf) accessible by both the CPU and the FPGA, and are never
re-transferred after the first use. The current activation vector is quantized from
FP32 to INT8 on the CPU. The CPU programs the accelerator via memory-mapped registers,
asserts a start signal, and sleeps until an interrupt from the FPGA fabric signals
completion. The FPGA computes the full MatMul + SwiGLU block in a pipelined dataflow
and writes the FP32 result back to the DMA buffer, from which the CPU copies it into
the GGML output tensor. All other layer operations remain on the CPU. If conditions
are not met at graph-build time or \texttt{LLAMA\_SWIHW} is set to 0, execution falls
back to the standard CPU path.

\begin{algorithm}[htb!]
\caption{LFM2.5 Inference Execution}
\label{alg:inference}
\begin{algorithmic}[1]
\STATE load\_model(GGUF)
\FOR{each token}
    \STATE $graph \leftarrow$ build\_ggml\_graph(token)
    \FOR{each node in $graph$}
        \IF{node.op == SWIGLU\_FUSED\_HW and enabled}
            \STATE // CPU--FPGA handoff
            \IF{not $weights\_cached[layer]$}
                \STATE Copy $W_{gate}, W_{up}, W_{down} \rightarrow shared\_mem$
                \STATE Mark cached
            \ENDIF
            \STATE $x_{int8}, x_{scale} =$ quantize($x$)
            \STATE Write $x_{int8} \rightarrow shared\_mem$
            \STATE program\_fpga($W_{addrs}, x_{addr}, out_{addr}, mode, x_{scale}$) \COMMENT{write AXI-Lite ctrl regs}
            \STATE start\_fpga()
            \STATE wait\_for\_interrupt()
            \STATE $out =$ read\_output($shared\_mem$)
        \ELSE
            \STATE $out =$ cpu\_kernel(node)
        \ENDIF
    \ENDFOR
    \STATE \COMMENT{graph execution complete; proceed to next token}
\ENDFOR
\end{algorithmic}
\end{algorithm}

Although Algorithm~\ref{alg:inference} expresses the offload as a single conditional
inside the per-node loop, its effect over a full token is dictated by the block
structure of Section~\ref{subsec:lfm2_bg}. The computation graph of one token contains
sixteen SwiGLU FFN nodes, one per transformer block, and the driver intercepts every
one of them. The offload is therefore not a single event but a sequence of sixteen
separate, synchronous CPU--FPGA handoffs per generated token.

These handoffs cannot be fused or chained on the FPGA, because consecutive FFNs are
never adjacent in the graph. The output of the FFN in block $N$ is consumed by the
residual addition, the RMSNorm, and the sequence mixer of block $N{+}1$, all of which
execute on the CPU; only their result becomes the input activation of the FFN in block
$N{+}1$. The accelerator thus has no way to advance from one FFN to the next on its own:
the CPU lies on the critical path between every pair of feed-forward blocks. As a
consequence the two devices strictly alternate: while the FPGA evaluates one FFN the CPU
is idle waiting on the completion interrupt, and while the CPU computes a block's mixer
and normalizations the FPGA is idle, so neither device overlaps the other within a token.
Figure~\ref{fig:handoff_mechanism} details the protocol of a single handoff.

\begin{figure}[htb!]
\centering
\begin{tikzpicture}[
  font=\scriptsize, >=Latex,
  head/.style={draw, rounded corners, fill=gray!12, minimum width=2.4cm, minimum height=8mm, align=center},
  cpuhead/.style={head, fill=blue!10},
  fpgahead/.style={head, fill=red!18, draw=red!60!black},
  life/.style={dashed, gray},
  msg/.style={-Latex, thick},
  comp/.style={draw, fill=red!10, align=center, minimum width=2.4cm, minimum height=8mm}
]
\def\xc{0}\def\xd{4.6}\def\xf{9.2}
% lifeline headers
\node[cpuhead]  (cpu)  at (\xc,0) {\textbf{CPU} (llama.cpp)};
\node[head]     (ddr)  at (\xd,0) {\textbf{udmabuf} (shared DDR)};
\node[fpgahead] (fpga) at (\xf,0) {\textbf{SwiGLU IP} (FPGA)};
% lifelines
\draw[life] (cpu)  -- (\xc,-6.3);
\draw[life] (ddr)  -- (\xd,-6.3);
\draw[life] (fpga) -- (\xf,-6.3);
% CPU blocked activation bar
\fill[blue!10, draw=blue!40] (\xc-0.12,-2.1) rectangle (\xc+0.12,-5.0);
\node[rotate=90, font=\tiny, text=blue!50!black] at (\xc-0.5,-3.55) {blocked};
% messages
\draw[msg] (\xc,-0.9) -- node[above]{(1) write $x$ (INT8)} (\xd,-0.9);
\draw[msg] (\xc,-1.5) -- node[above]{(2) program control registers (AXI-Lite)} (\xf,-1.5);
\draw[msg] (\xc,-2.0) -- node[above]{(3) assert \texttt{ap\_start}} (\xf,-2.0);
\draw[msg] (\xf,-2.7) -- node[above]{(4) read $x$ and cached weights (AXI4)} (\xd,-2.7);
\node[comp] at (\xf,-3.3) {(5) compute 5-stage FFN};
\draw[msg] (\xf,-4.3) -- node[above]{(6) write out (FP32)   } (\xd,-4.3);
\draw[msg] (\xf,-5.0) -- node[above]{(7) interrupt (\texttt{ap\_done})} (\xc,-5.0);
\draw[msg] (\xc,-5.7) -- node[above]{(8) read out} (\xd,-5.7);
% execution-flow indicator (time runs downward)
\draw[-Latex, thick, gray] (-1.8,-0.6) -- (-1.8,-6.0);
\node[rotate=90, font={\scriptsize\itshape}, text=gray] at (-2.1,-3.3) {execution flow};
\end{tikzpicture}
\caption{Protocol of a single CPU--FPGA handoff. The CPU quantizes the activation and
writes it to the shared udmabuf buffer, programs the IP's AXI-Lite control registers,
asserts \texttt{ap\_start} and blocks; the FPGA burst-reads the activation and the cached
weights, evaluates the five-stage SwiGLU pipeline, writes the result back, and raises a
completion interrupt that releases the CPU to copy the output. This sequence repeats
sixteen times per generated token.}
\label{fig:handoff_mechanism}
\end{figure}

Each of the sixteen handoffs carries a fixed overhead that is independent of the work it
guards: the input activation is quantized to INT8, the AXI-Lite control registers are
programmed, the shared buffer is synchronized for the device, the calling thread blocks
until the completion interrupt, and the result is copied back into the GGML tensor. This
cost is paid sixteen times per token and does not amortize with sequence length the way
the arithmetic of a batched matrix multiply would, since decode generates one token at a
time. It is the principal component of the $\approx 126$\,ms of non-FFN per-token
overhead measured in Section~\ref{subsec:cpuexec}, which fixes the Amdahl ceiling near
8 tokens per second. The synchronous, blocking nature of the handoff also explains the
behaviour at high thread counts (Section~\ref{sec:discussion}): when all four cores are
busy, the blocking wait competes with the interrupt handler for scheduler time and the
per-token round-trip latency grows, so the FPGA path can fall below the CPU-only
baseline. Overlapping FPGA execution with independent CPU work would require an
asynchronous, pipelined handoff in place of the node-by-node, barrier-synchronized
execution used here; this direction, and the compatibility trade-off it entails, is
discussed in Chapter~\ref{cha:conclusioni}.

\section{FPGA Accelerator}
\label{sec:fpga}

The proposed accelerator implements the SwiGLU feed-forward block of the LFM2.5 language
backbone,
\begin{equation}
\text{FFN}(x) = \big(\text{SiLU}(x W^{T}) \odot (x V^{T})\big)\, W_{down}^{T},
\label{eq:ffn_hw}
\end{equation}
which the five-stage dataflow pipeline evaluates through the decomposition:
\begin{subequations}\label{eq:swiglu_stage_set}
\begin{align}
X_1 &= x W^{T} \label{eq:swiglu_s2a}\\
X_2 &= x V^{T} \label{eq:swiglu_s2b}\\
G   &= \text{SiLU}(X_1)\odot X_2 \label{eq:swiglu_s3}\\
\text{FFN}(x) &= G W_{down}^{T} \label{eq:swiglu_s4}
\end{align}
\end{subequations}
where $x$ is the input activation, $W$ is the gating projection, $V$ is the expansion
projection, and $W_{down}$ is the output projection; $\odot$ is element-wise
multiplication and $\text{SiLU}(z) = z \cdot \sigma(z)$. Equation~\eqref{eq:swiglu_s2a}
is computed in Stage~2a (gate projection), Eq.~\eqref{eq:swiglu_s2b} in Stage~2b (up
projection), Eq.~\eqref{eq:swiglu_s3} in Stages~3--4 (gated activation: SiLU sub-stage
A, element-wise multiply sub-stage B), and Eq.~\eqref{eq:swiglu_s4} in Stage~5 (down
projection).

The design exposes AXI4 interfaces for weights and activations, and an AXI-Lite
interface for runtime control (quantization mode and per-token scaling factors). The
five-stage pipeline is organized as an HLS dataflow (Fig.~\ref{fig:swiglu_compact})
targeting an initiation interval $II=1$ throughout; end-to-end latency is therefore
determined by memory bandwidth and burst efficiency rather than arithmetic capacity.

\begin{figure}[htb!]
\centering
\begin{tikzpicture}[
    font=\scriptsize,
    >=Latex,
    node distance=5mm,
    block/.style={draw, rounded corners, thick, align=center, minimum height=6mm, text width=2.6cm},
    stageWorst/.style={block, fill=red!18, draw=red!60!black},
    stageLow/.style={block, fill=yellow!40, draw=yellow!60!black},
    stageStd/.style={block, fill=gray!15, draw=gray!60},
    mem/.style={draw, thick, align=center, minimum height=5mm, text width=2.4cm, fill=gray!10},
    line/.style={-Latex, thick}
]

\node[stageStd] (loadx) {\textbf{Stage 1}\\Load $x$ (INT8)\\202 cy / 0.81\,\textmu s};

\node[mem, below left=3mm and 2mm of loadx] (x1buf) {$x_1$ bank};
\node[mem, below right=3mm and 2mm of loadx] (x2buf) {$x_2$ bank};

\node[stageLow, below=7mm of x1buf] (x1) {\textbf{Stage 2a}\\$X_1=xW_{gate}^{T}$\\1.54 Million cycles / 6.14\,ms};
\node[stageLow, below=7mm of x2buf] (x2) {\textbf{Stage 2b}\\$X_2=xW_{up}^{T}$\\1.54 Million cycles / 6.14\,ms};

\node[mem, below=of x1] (x1cache) {$X_1$};
\node[mem, below=of x2] (x2cache) {$X_2$};

\node[stageStd, below=7mm of x1cache] (silu) {\textbf{Stage 3} {\scriptsize(sub-stage A)}\\SiLU($X_1$)\\{\scriptsize$\approx$8.0K cy / 32\,\textmu s}};

\node[stageStd, below=10mm of $(silu)!0.5!(x2cache)$] (gate)
{\textbf{Stage 4} {\scriptsize(sub-stage B)}\\gate=SiLU($X_1$)$\cdot X_2$\\{\scriptsize$\approx$8.6K cy / 34\,\textmu s}};

\node[stageWorst, below=of gate] (down)
{\textbf{Stage 5}\\$gate \cdot W_{down}^{T}$\\1.62 Million cycles / 6.49\,ms};

\node[mem, below=of down] (out) {out};

\draw[line] (loadx) -- (x1buf);
\draw[line] (loadx) -- (x2buf);
\draw[line] (x1buf) -- (x1);
\draw[line] (x2buf) -- (x2);
\draw[line] (x1) -- (x1cache);
\draw[line] (x2) -- (x2cache);
\draw[line] (x1cache) -- (silu);
\draw[line] (silu.south) |- (gate.west);
\draw[line] (x2cache.south) |- (gate.east);
\draw[line] (gate) -- (down);
\draw[line] (down) -- (out);

\node[draw, rounded corners, fill=white, align=left, anchor=north, font=\scriptsize]
at ($(out.south)+(0mm,-6mm)$) {
\fcolorbox{gray!60}{gray!15}{\strut\hspace{0.9em}}~Negligible DDR traffic, not load-bound\\
\fcolorbox{yellow!60!black}{yellow!40}{\strut\hspace{0.9em}}~$\geq$40\% DDR bandwidth utilization, load-bound\\
\fcolorbox{red!75!black}{red!18}{\strut\hspace{0.9em}}~$<$40\% DDR bandwidth utilization, load-bound, DATAFLOW bottleneck};

\end{tikzpicture}
\caption{Compact dataflow of the MatMul + SwiGLU accelerator.}
\label{fig:swiglu_compact}
\end{figure}

\subsection{HLS Realization}
\label{subsec:hls}

The accelerator is written in C++ and compiled to register-transfer-level hardware by
Vitis HLS. Its functional behaviour is expressed in ordinary loops and arithmetic, while
the mapping to hardware (interface protocols, on-chip memories, parallelism, and
pipelining) is steered by directives, or pragmas, that annotate the code without altering
its result. Listing~\ref{lst:hls} shows the top-level function, which both defines the
IP's external interface and launches the five pipeline stages.

\begin{lstlisting}[style=hls, float=htb, caption={Condensed top-level HLS function: the external interface and the five-stage dataflow. Repeated \texttt{m\_axi} and \texttt{s\_axilite} lines are collapsed for clarity.}, label=lst:hls]
void swiglu(const uint8_t *W, const uint8_t *V, const uint8_t *W_down,
            const int8_t *x_batch, float *out_batch,
            uint32_t down_quant_mode, float x_scale) {
  // weight/activation pointers become 128-bit AXI4 master ports to DDR
  #pragma HLS INTERFACE m_axi port=W bundle=gmem_W max_widen_bitwidth=128
  // ... V, W_down, x_batch, out_batch likewise ...
  // scalar arguments become memory-mapped control registers
  #pragma HLS INTERFACE s_axilite port=x_scale bundle=CTRL
  // ... addresses, down_quant_mode, ap_ctrl likewise ...

  int8_t X1_cache[FFN_DIM];
  #pragma HLS BIND_STORAGE    variable=X1_cache   type=ram_2p impl=bram
  int8_t gate_cache[DOWN_BLOCKS][256];
  #pragma HLS BIND_STORAGE    variable=gate_cache type=ram_1p impl=uram
  #pragma HLS ARRAY_PARTITION variable=gate_cache dim=2 cyclic factor=8

  #pragma HLS DATAFLOW                          // stages run as overlapping processes
  load_x_local(x_batch, x_local_1, x_local_2);
  compute_X1(W, x_local_1, x_scale, X1_cache);   // Stage 2a
  compute_X2(V, x_local_2, x_scale, X2_cache);   // Stage 2b
  compute_gate(X1_cache, X2_cache, gate_cache);  // Stages 3-4
  compute_output(W_down, gate_cache, out_batch); // Stage 5
}
\end{lstlisting}

Four kinds of directive appear. The \texttt{m\_axi} pragmas turn the weight and
activation pointers into AXI4 master ports. Since HLS by default infers the bus width
from the \texttt{uint8\_t} element type, giving an 8-bit bus, the
\texttt{max\_widen\_bitwidth=128} option is what widens each port to 128 bits, the width
of the device's PS--PL high-performance interface, and raises its achievable DDR
bandwidth roughly sixteen-fold relative to the 8-bit default. Widening the port further brings no benefit, since it cannot exceed that 128-bit interface width. 
The \texttt{s\_axilite} pragmas expose the
scalar arguments (the weight and buffer addresses, the quantization mode, the per-token
scale, and the \texttt{ap\_ctrl} start/done handshake) as the memory-mapped control
registers that the driver programs during the handoff of
Section~\ref{sec:orchestration}. The \texttt{BIND\_STORAGE} directives pin each on-chip
array to a specific primitive, LUTRAM, BRAM, or URAM, so that the design fits the
ZU5EV's mixed memory budget; \texttt{X1\_cache} is a dual-port (\texttt{ram\_2p}) BRAM
because it is written by \texttt{compute\_X1} and read by \texttt{compute\_gate}
concurrently. \texttt{ARRAY\_PARTITION} with \texttt{cyclic factor=8} splits the gate
cache into eight URAM banks so that eight blocks can be read in a single cycle, feeding
the parallel multiply-accumulate chains. A larger factor would only add banks that the
multiply-accumulate stage never reads in the same cycle; eight matches the eight-block
width of the rest of the pipeline. Finally, \texttt{DATAFLOW} instructs the tool to
run the five stage functions as concurrent, overlapping processes rather than in
sequence: this is what realizes the pipeline of Figure~\ref{fig:swiglu_compact}, with the
two projection stages executing in parallel.

The remaining directives operate inside the stage functions: each accumulation loop is
pipelined at an initiation interval of one (\texttt{PIPELINE II=1}) with the block loop
unrolled (\texttt{UNROLL}), so the partitioned weight and activation banks are consumed in
parallel. This is the mechanism behind the 32 multiply-accumulate chains detailed in the
following subsections.

\subsection{Quantization}

Q4\_K is preferred over the structurally simpler Q4\_0 format (one fp16 block scale per
32 elements, no sub-block scales) on quality grounds: its 6-bit sub-group scales
provide $16\times$ finer scale resolution per block, keeping quantization fidelity
above the Q4\_0 baseline that Liquid AI themselves use for all LFM2 inference
benchmarks~\cite{amini2025lfm2technicalreport}. The default GGUF distribution of
LFM2.5-1.2B applies mixed quantization to the down-projection weights, using Q6\_K for
eight layers and Q4\_K for the remaining eight (the Q4\_K\_M strategy). Supporting
Q6\_K in the accelerator is infeasible on the ZU5EV: each Q6\_K block occupies 210
bytes ($13.125\times$128-bit words), precluding the clean 2D register-file layout that
enables compile-time \texttt{.range()} nibble extraction. The fallback
\texttt{get\_byte()} multiplexer tree required for 6-bit nibbles consumes approximately
30--40\,K LUTs per MAC stage at 32-block parallelism, exceeding the device's remaining
routing capacity at 87\% LUT occupancy. The model is therefore re-quantized to uniform
Q4\_K across all sixteen layers using
\texttt{llama-quantize -{-}allow-requantize -{-}pure} from an F16\_Q8\_0 source. This
choice is consistent with the hardware-in-the-loop co-design process of LFM2, whose
architecture search and published inference benchmarks evaluated all candidates under
Q4\_0 quantization~\cite{amini2025lfm2technicalreport}, a scheme coarser than Q4\_K in
both bit-width and block-scale granularity. Furthermore, Q6\_K is applied exclusively
to the down-projection $W_{down}$, the least quantization-sensitive weight in the FFN:
it acts as a linear readout of an already-nonlinear gated intermediate, and its
32-block structure averages quantization error before injection into the residual
stream. Across the full mixed-precision datapath, accuracy degradation relative to the
CPU's FP32 reference is estimated from quantization error bounds: the dominant
contribution is Q4\_K weight error (approximately 6\% relative to the FP32 dynamic
range), with INT8 input quantization adding under 0.5\% and INT8 intermediate cache
quantization contributing at most 1\%. These are order-of-magnitude bounds, not direct
measurements on this model; the impact on task-level accuracy is expected to be within
the tolerance of a conversational audio deployment.

\subsection{CPU-Side Weight Pre-Decoding}
\label{subsec:predecode}

The $K=4$ designation indicates that four weight rows are processed per iteration of each
MAC stage, yielding 32 parallel INT32 multiply-accumulate chains (8 blocks $\times$ 4
rows) in the projection stages and 32 chains shared across 4 sequential groups in the
down-projection stage. Sustaining this degree of parallelism within the device's LUT
budget is made possible by a pre-decoding step that runs on the CPU, once per layer,
before the weights are ever seen by the accelerator.

The difficulty stems from the Q4\_K storage format, which is designed for CPU SIMD
execution rather than for an FPGA datapath. Each 144-byte block encodes 256 weights as
shown in Table~\ref{tab:q4k_block}: two fp16 block-level scales ($d$ and $dmin$), eight
6-bit sub-block scales and eight 6-bit minima (\texttt{sc6}, \texttt{mn6}) packed at
bit-level offsets that straddle byte boundaries, and 256 four-bit nibbles stored in
block-major order. Decoding this layout inside the MAC pipeline would require, for every
nibble, a \texttt{get\_byte()} access, a nine-to-one multiplexer on the 128-bit word
followed by a four-bit shift. With eight blocks decoded in parallel, these multiplexer
trees alone consume on the order of 30\,K LUTs, which is incompatible with the $K=4$
design on the ZU5EV.

\begin{table}[htbp]
\centering
\caption{Layout of a 144-byte Q4\_K block (256 weights).}
\label{tab:q4k_block}
\begin{tabular}{lll}
\hline
\textbf{Bytes} & \textbf{Size} & \textbf{Field} \\ \hline
0--1    & 2\,B   & $d$ (fp16 block scale) \\
2--3    & 2\,B   & $dmin$ (fp16 min-subtraction scale) \\
4--11   & 8\,B   & \texttt{sc6}/\texttt{mn6} sub-block scales and minima (low 6 bits) \\
12--15  & 4\,B   & \texttt{sc6}/\texttt{mn6} high-order bits (interleaved) \\
16--143 & 128\,B & 256 packed 4-bit nibbles (block-major) \\ \hline
\end{tabular}
\end{table}

The pre-decoding step performs two transformations that move all of this irregular bit
manipulation off the critical path and onto the CPU. First, the interleaved 6-bit
\texttt{sc6} and \texttt{mn6} fields are unpacked into flat INT8 arrays, so the FPGA reads
each sub-block scale as a plain byte. Second, and more importantly, the nibbles are
transposed from the block-major order in which Q4\_K stores them into an element-major
order, illustrated in Figure~\ref{fig:predecode}: instead of all 256 nibbles of block~0
followed by all 256 of block~1, the pre-decoded layout places the nibble of a given
element for all eight blocks side by side in a single 128-bit word. The fp16 scales $d$
and $dmin$ are copied verbatim and converted to fp32 in hardware by a custom
bit-manipulation routine that preserves subnormal values rather than flushing them to
zero.

\begin{figure}[htb!]
\centering
\begin{tikzpicture}[
  font=\scriptsize, >=Latex,
  raw/.style={draw, rounded corners, fill=gray!12, align=center, text width=4.0cm, minimum height=14mm},
  pre/.style={draw, rounded corners, fill=blue!10, align=center, text width=4.4cm, minimum height=14mm}
]
\node[raw] (a) {Raw Q4\_K (block-major):\\all 256 nibbles of block~0,\\then block~1, \dots, block~7};
\node[pre, right=2.6cm of a] (b) {Pre-decoded (element-major):\\one 128-bit word per element\\holds that element's nibble\\for all 8 blocks ($b_0\dots b_7$)};
\draw[-Latex, thick] (a) -- node[above, align=center]{CPU pre-decode\\(once / layer)} (b);
\end{tikzpicture}
\caption{Weight pre-decoding rearranges the block-major Q4\_K nibbles into an
element-major layout, so each 128-bit memory word holds one element's nibble for all
eight blocks and can be consumed by the parallel MAC chains in a single cycle.}
\label{fig:predecode}
\end{figure}

After this rearrangement the FPGA load is trivial: each 128-bit DDR word is written
verbatim into a wide (128-bit) BRAM tile, one read and one write per cycle, achieving an
initiation interval of one with no decode logic at all. In the MAC, the eight nibbles of a
word are addressed by compile-time \texttt{.range()} slices, which synthesize to wires and
cost zero LUTs. Eliminating the multiplexer trees is what brings the $K=4$ implementation
from roughly 137\,K LUTs (which does not fit the ZU5EV's 117\,K) down to roughly
103\,K, leaving margin for routing; the implemented design occupies 87\% of the LUTs
(Section~\ref{sec:results}). The pre-decode itself costs about
10\,ms per layer on the first token; its output is cached permanently in the shared
udmabuf buffer, so every subsequent token reads the transposed weights directly and pays
nothing.

\subsection{Stage 1: Load $x$ (INT8)}
The 2048-element INT8 input vector is burst-loaded from DDR in $128\times$128-bit AXI
words and duplicated into two identically partitioned LUTRAM buffers, enabling
concurrent access by Stages~2a and~2b without serialization. The stage is pure data
movement with no arithmetic and completes in 202 cycles (0.81\,$\mu$s).

\subsection{Stages 2a / 2b: Linear Projections (Parallel)}
Two GEMV operations compute $X_1$ and $X_2$
(Eqs.~\eqref{eq:swiglu_s2a}--\eqref{eq:swiglu_s2b}) in parallel on independent AXI ports
and BRAM banks. Weights are pre-decoded on the CPU into flat INT8 sub-block scales and
element-major nibble words; the FPGA never handles the interleaved Q4\_K header format.
With $K=4$, four rows are processed per outer iteration via a single merged 320-word AXI
burst covering all four rows simultaneously (16 header + 64 nibble words per row,
$\times\,4$ rows = 320 total), eliminating the per-iteration burst-setup overheads of a
design with multiple calls. Nibble extraction uses compile-time \texttt{.range()}
indexing on wide 128-bit BRAM tiles. The 32 parallel INT32 MAC chains (8 blocks $\times$
4 rows) deliver 2.76\,GOPS (34.5\% of the 8.00\,GOPS peak); the merged burst reduces
load time from 948 to 401 cycles per iteration, representing 53\% of the total
per-iteration cycle budget. DDR bandwidth reaches 1.63\,GB/s (41\% of the 4.00\,GB/s
peak) across the 10.0\,MB per-matrix transfer. Each stage completes in 1,536,001 cycles
(6.14\,ms).

\subsection{Stages 3--4: Gated Activation and Quantization}
The fused gate operation (Eq.~\eqref{eq:swiglu_s3}) runs entirely on-chip from the
$X_1$ and $X_2$ caches using a two-pass scheme over 8,192 elements: the first pass finds
the maximum absolute value across 8 parallel reduction lanes, and the second recomputes
and requantizes to INT8 for storage in the 8-way banked URAM gate cache. At 16,622
cycles (66\,$\mu$s) and approximately 65K FP operations per pass, the stage contributes
under 1\% of per-layer time.

\subsection{Stage 5: Down Projection}
The gated activation is projected back to model dimension (Eq.~\eqref{eq:swiglu_s4}) via
$\text{out} = \text{gate} \cdot W_{down}^{T}$ using the same CPU-pre-decoded weight
format as Stages~2a/2b. With $K=4$, four output rows are processed per iteration across
32 blocks organized into 4 sequential groups of 8; all four rows are unrolled per group
to deliver 32 parallel INT32 MAC chains per group. Unlike Stages~2a/2b, which merge all
four rows into a single 320-word burst, Stage~5 issues four sequential per-row AXI bursts
of 320 words each. The merged strategy is not applicable here: $W_{down}$ has 32 blocks
per row (versus 8 in $W$/$V$), so each row's nibble buffer is 256 words rather than 64.
Consequently, nibbles are stored in narrower 32-bit tiles (4 groups $\times$ 4 rows)
rather than 128-bit tiles. This tile layout also precludes the contiguous multi-row
memory mapping that enables burst merging in Stages~2a/2b. These loads account for 60\%
of the per-iteration budget (1,908 of 3,165 cycles), making the stage load-bound. DDR
bandwidth reaches 1.54\,GB/s (38\% of the 4.00\,GB/s peak) across the 10.0\,MB transfer.
Total latency is 1,622,602 cycles (6.49\,ms), which sets the dataflow pipeline interval.
